\chapter{The Symmetric Group}

\begin{center}\itshape
The irreducible representations of $\mathfrak S_n$ are indexed by the partitions of $n$, and each one can be built by hand from a Young diagram. Its dimension is the number of ways to fill the diagram with $1,\dots,n$ increasingly along rows and
columns, and a formula with hooks counts these fillings. A bijection between permutations and pairs of fillings proves it, and the same bijection measures the longest increasing subsequence of a permutation.
\end{center}

Alfred Young was born in Widnes, Lancashire, in 1873, was taught at home and then at Monkton Combe school, and came tenth Wrangler at Cambridge in 1895. He worked on the invariant theory of binary forms, the
kind of symbolic algebra that the Cambridge of his day loved, and in a paper of 1900 he started to draw rows of boxes filled with numbers. In 1901 and 1902 the first two papers of a series called \emph{On quantitative
substitutional analysis} appeared; ``substitutions'' were permutations, and the boxes, which we now call \emph{Young tableaux}, were his tool for decomposing the group algebra of $\mathfrak S_n$.
Herbert Turnbull later wrote: ``From the outset Young's rapid and skilful handling of symbolic algebra bore all the signs of genius.''

Then Young went quiet. He became a fellow and bursar of Clare College, was ordained in 1908, and in 1910 took the parish of Birdbrook, a village in Essex, where he lived for the rest of his life. Meanwhile in Berlin, Frobenius computed the
characters of $\mathfrak S_n$ (1900) and used Young's diagrams (1903); Young, who did not read German easily, only learned of this later, and he took up his series again in the late 1920s. Nine papers appeared in all, the last one in 1952, eleven years after
his death. In 1934 the country parson was elected a Fellow of the Royal Society.

This chapter tells the modern version of Young's story. Section~1 builds, for each partition $\lambda$ of $n$, a representation $S^\lambda$ (Specht, 1935), and proves with a short argument of Gordon James (1978) that these are all the
irreducible representations of $\mathfrak S_n$. Section~2 finds a basis indexed by \emph{standard tableaux}. Section~3 counts them twice: with the Robinson--Schensted correspondence, a bijection between permutations and pairs of tableaux,
and with the hook length formula of Frame, Robinson and Thrall (1954), proved by a random walk inside the diagram. Section~4 computes the character of a transposition, which is exactly what Volume~11 uses to shuffle cards.

Throughout, $n\ge1$, $\mathfrak S_n$ acts on $\{1,\dots,n\}$, and representations are complex and finite-dimensional, as in Volume~3, Chapter~11. Permutations are composed right to left: $\sigma\pi$ means $\pi$ first.

\section{Young diagrams and Specht modules}

\begin{definition}\label{def:partition}
A \emph{partition} $\lambda\vdash n$ is a sequence $\lambda_1\ge\lambda_2\ge\dots\ge\lambda_k>0$ of integers with sum $n$. Its \emph{Young diagram} is the set of cells $(i,j)$ with $1\le i\le k$ and $1\le j\le\lambda_i$, drawn as boxes, row $i$ from the top,
column $j$ from the left. The \emph{conjugate} partition $\lambda'$ lists the column lengths: $\lambda'_j=\#\{i:\lambda_i\ge j\}$. A \emph{tableau} of shape $\lambda$ is a filling $T$ of the cells with $1,\dots,n$, each used once.
\end{definition}

The conjugacy classes of $\mathfrak S_n$ are the cycle types (Volume~3, Chapter~4), and a cycle type is a partition of $n$: so $\mathfrak S_n$ has as many irreducible representations as there are partitions of $n$ (Volume~3, Theorem~11.10).
The point of this section is that there is a natural bijection, not just an equality of numbers.

A permutation $\sigma$ acts on tableaux by replacing each entry $a$ with $\sigma(a)$; write $\sigma T$. The \emph{row group} $R_T$ is the subgroup of permutations that preserve each row of $T$ as a set, and the \emph{column group}
$C_T$ the one that preserves each column. Then $R_{\sigma T}=\sigma R_T\sigma^{-1}$ and $C_{\sigma T}=\sigma C_T\sigma^{-1}$, and $R_T\cap C_T=\{1\}$: an entry is the only one in both its row and its column.

\begin{definition}\label{def:tabloid}
Two tableaux $T,T'$ of shape $\lambda$ are \emph{row equivalent} if $T'=\pi T$ for some $\pi\in R_T$, that is, if they have the same rows as sets. The class $\{T\}$ is a \emph{tabloid}. The permutation module $M^\lambda$ is the
vector space with basis the tabloids of shape $\lambda$, with $\sigma\{T\}=\{\sigma T\}$. For a tableau $T$, let
\[
\kappa_T=\sum_{\pi\in C_T}\operatorname{sgn}(\pi)\,\pi\in\C[\mathfrak S_n],\qquad e_T=\kappa_T\{T\}\in M^\lambda.
\]
The element $e_T$ is a \emph{polytabloid}, and the \emph{Specht module} $S^\lambda$ is the subspace of $M^\lambda$ spanned by all polytabloids.
\end{definition}

The action on tabloids is well defined, since $\sigma$ maps the rows of $T$ to the rows of $\sigma T$. It is transitive, and the stabiliser of $\{T\}$ is $R_T\cong\mathfrak S_{\lambda_1}\times\dots\times\mathfrak S_{\lambda_k}$, the
\emph{Young subgroup} $\mathfrak S_\lambda$; so $M^\lambda$ is the permutation representation on $\mathfrak S_n/\mathfrak S_\lambda$, which is $\operatorname{Ind}_{\mathfrak S_\lambda}^{\mathfrak S_n}1$ (Chapter~9), of dimension $n!/(\lambda_1!\cdots\lambda_k!)$.
Since $\kappa_{\sigma T}=\sigma\kappa_T\sigma^{-1}$, we have $\sigma e_T=e_{\sigma T}$: the span $S^\lambda$ is a submodule, generated by any single polytabloid.

For example, $M^{(n)}$ is the trivial representation, and $S^{(n)}=M^{(n)}$. For $\lambda=(1^n)$, one cell per row, tabloids are tableaux, $M^{(1^n)}=\C[\mathfrak S_n]$, and $C_T=\mathfrak S_n$; then $\sigma e_T=e_{\sigma T}=\operatorname{sgn}(\sigma)e_T$ (substitute
$\pi\mapsto\pi\sigma^{-1}$ in the sum), so $S^{(1^n)}$ is the sign representation. For $\lambda=(n-1,1)$, a tabloid is determined by the entry $a$ of its second row, so $M^{(n-1,1)}=\C^n$ with basis vectors $\mathbf e_a$; if the first column of $T$ holds $b$
above $a$, then $e_T=\mathbf e_a-\mathbf e_b$, and $S^{(n-1,1)}$ is the hyperplane $\sum x_a=0$, the standard representation of dimension $n-1$.

\begin{definition}\label{def:dominance}
For partitions $\lambda,\mu$ of $n$, $\lambda$ \emph{dominates} $\mu$, written $\lambda\trianglerighteq\mu$, if $\lambda_1+\dots+\lambda_i\ge\mu_1+\dots+\mu_i$ for every $i$ (with $\lambda_i=0$ beyond the last part).
\end{definition}

Dominance is a partial order: $(3,1)\trianglerighteq(2,2)\trianglerighteq(2,1,1)$, but $(3,1,1,1)$ and $(2,2,2)$ are not comparable. The next lemma is the engine of the chapter.

\begin{lemma}[Dominance lemma]\label{lem:dominance}
Let $T$ be a tableau of shape $\lambda$ and $T'$ one of shape $\mu$. If, for every $i$, the entries of row $i$ of $T'$ lie in different columns of $T$, then $\lambda\trianglerighteq\mu$.
\end{lemma}

\begin{proof}
Fix $i$ and look at the $\mu_1+\dots+\mu_i$ entries in the first $i$ rows of $T'$. Two entries of a column of $T$ come from different rows of $T'$, so a column of $T$ contains at most $i$ of them, and at most its length $\lambda'_j$.
Hence
\[
\mu_1+\dots+\mu_i\le\sum_j\min(i,\lambda'_j)=\lambda_1+\dots+\lambda_i,
\]
where the last equality counts the cells in the first $i$ rows of the diagram of $\lambda$ column by column.
\end{proof}

\begin{lemma}\label{lem:kappa}
Let $T$ be a tableau of shape $\lambda$ and $T'$ one of shape $\mu$.
\begin{enumerate}[label=(\alph*)]
\item If $\kappa_T\{T'\}\ne0$, then $\lambda\trianglerighteq\mu$.
\item If $\lambda=\mu$, then $\kappa_T\{T'\}\in\{0,e_T,-e_T\}$; consequently, $\kappa_Tu$ is a multiple of $e_T$ for every $u\in M^\lambda$.
\end{enumerate}
\end{lemma}

\begin{proof}
Suppose two entries $a,b$ lie in the same row of $T'$ and the same column of $T$. The transposition $t=(a\,b)$ is in $C_T$, and $t\{T'\}=\{T'\}$. Choose representatives $\pi_1,\dots,\pi_m$ of the left cosets of $\{1,t\}$ in $C_T$: then
$C_T=\{\pi_1,\pi_1t,\dots,\pi_m,\pi_mt\}$ and
\[
\kappa_T=\Bigl(\sum_{l}\operatorname{sgn}(\pi_l)\pi_l\Bigr)(1-t),\qquad\text{so}\qquad\kappa_T\{T'\}=\Bigl(\sum_{l}\operatorname{sgn}(\pi_l)\pi_l\Bigr)\bigl(\{T'\}-\{T'\}\bigr)=0.
\]
So if $\kappa_T\{T'\}\ne0$, the hypothesis of the dominance lemma holds, and $\lambda\trianglerighteq\mu$: this is (a).

For (b), assume moreover $\lambda=\mu$. We show that $\{T'\}=\pi\{T\}$ for some $\pi\in C_T$. The $\lambda_1$ entries of row $1$ of $T'$ lie in different columns of $T$, and $T$ has $\lambda_1$ columns, so each column of $T$ contains exactly one
of them; a permutation in $C_T$ moves each one to the top of its column, and the new tableau $T_1=\pi_1T$ has the same first row as $T'$. Now repeat with rows $2,3,\dots$ of $T'$ and the part of $T_1$ below row $1$, using
column permutations that fix row $1$. At the end, $\pi T$ has the same rows as $T'$ for some $\pi\in C_T$. Then $\kappa_T\pi=\operatorname{sgn}(\pi)\kappa_T$ (substitute $\pi'\mapsto\pi'\pi^{-1}$ in the sum), so
$\kappa_T\{T'\}=\kappa_T\pi\{T\}=\operatorname{sgn}(\pi)e_T$. The last claim follows by linearity, as tabloids span $M^\lambda$.
\end{proof}

Give $M^\lambda$ the Hermitian inner product for which the tabloids are orthonormal. Permutations permute tabloids, so they are unitary, and the adjoint of $\pi$ is $\pi^{-1}$; as $\operatorname{sgn}(\pi^{-1})=\operatorname{sgn}(\pi)$ and
$C_T$ is a group, $\kappa_T$ is self-adjoint.

\begin{theorem}[James' submodule theorem, 1978]\label{thm:james}
Let $U$ be a submodule of $M^\lambda$. Then either $U\supseteq S^\lambda$ or $U\subseteq(S^\lambda)^\perp$.
\end{theorem}

\begin{proof}
Let $u\in U$ and $T$ a tableau of shape $\lambda$. By Lemma~\ref{lem:kappa}(b), $\kappa_Tu=c\,e_T$ for a scalar $c$, and $\kappa_Tu\in U$. If $c\ne0$ for some $u$ and $T$, then $e_T\in U$, and $S^\lambda\subseteq U$ because $e_T$
generates $S^\lambda$. Otherwise $\kappa_Tu=0$ for all $u\in U$ and all $T$, and then
\[
\langle u,e_T\rangle=\langle u,\kappa_T\{T\}\rangle=\langle\kappa_Tu,\{T\}\rangle=0,
\]
so $U$ is orthogonal to every polytabloid, that is, $U\subseteq(S^\lambda)^\perp$.
\end{proof}

\begin{theorem}[Young, Frobenius, Specht]\label{thm:specht}
The Specht modules $S^\lambda$, for $\lambda\vdash n$, are irreducible and pairwise non-isomorphic, and every irreducible representation of $\mathfrak S_n$ is isomorphic to exactly one of them.
\end{theorem}

\begin{proof}
\emph{Irreducible.} $S^\lambda\ne0$, since $e_T$ contains $\{T\}$ with coefficient $1$: if $\pi\in C_T$ and $\pi\{T\}=\{T\}$, then $\pi\in C_T\cap R_T$, so $\pi=1$. Let $U\subseteq S^\lambda$ be a submodule. By Theorem~\ref{thm:james}, either $U=S^\lambda$
or $U\subseteq S^\lambda\cap(S^\lambda)^\perp=0$, the inner product being positive definite.

\emph{Non-isomorphic.} Suppose $\theta:S^\lambda\to M^\mu$ is a non-zero module map. The orthogonal complement $(S^\lambda)^\perp$ is a submodule (the action is unitary), so $M^\lambda=S^\lambda\oplus(S^\lambda)^\perp$, and $\theta$ extends to a module
map $\Theta:M^\lambda\to M^\mu$ that is zero on $(S^\lambda)^\perp$. Since $\theta\ne0$, some $e_T$ has $\theta(e_T)\ne0$, and
\[
\theta(e_T)=\Theta(\kappa_T\{T\})=\kappa_T\,\Theta(\{T\}).
\]
Writing $\Theta(\{T\})$ as a combination of tabloids of shape $\mu$, some $\kappa_T\{T'\}$ is non-zero, so $\lambda\trianglerighteq\mu$ by Lemma~\ref{lem:kappa}(a). If $S^\lambda\cong S^\mu$, there are non-zero maps $S^\lambda\to S^\mu\subseteq M^\mu$ and
$S^\mu\to S^\lambda\subseteq M^\lambda$, so $\lambda\trianglerighteq\mu$ and $\mu\trianglerighteq\lambda$, hence $\lambda=\mu$, dominance being antisymmetric.

\emph{Complete.} We have found as many pairwise non-isomorphic irreducible representations as there are partitions of $n$, which is the number of conjugacy classes, hence the number of irreducible representations.
\end{proof}

The proof gives more: every irreducible constituent $S^\lambda$ of $M^\mu$ satisfies $\lambda\trianglerighteq\mu$, and $S^\mu$ occurs in $M^\mu$ exactly once (Exercise~\ref{ex:youngmult}). The exact multiplicities, the \emph{Kostka numbers}, count
fillings of $\lambda$ with $\mu_1$ ones, $\mu_2$ twos, and so on, increasing weakly along rows and strictly down columns (Young's rule). We will not need them.

\begin{think}
Over a field of characteristic $p\le n$, the same constructions make sense, but the inner product is no longer definite: $S^\lambda\cap(S^\lambda)^\perp$ can be non-zero. Which step of the proof of Theorem~\ref{thm:specht} breaks? James showed that
the quotients $S^\lambda/(S^\lambda\cap(S^\lambda)^\perp)$ are still irreducible or zero; which are which is a hard open problem for general $p$ and $n$.
\end{think}

\section{Standard tableaux}

\begin{definition}\label{def:standard}
A tableau is \emph{standard} if its entries increase along each row (left to right) and down each column. The number of standard tableaux of shape $\lambda$ is written $f^\lambda$.
\end{definition}

For $\lambda=(2,2)$ there are two standard tableaux, with rows $12/34$ and $13/24$. For $(n-1,1)$ there are $n-1$: the entry below $1$ can be any of $2,\dots,n$. In a standard tableau, the largest entry $n$ sits at the end of its row
and the bottom of its column, in a cell that can be removed leaving a partition: a \emph{corner}. Removing it gives a standard tableau of shape $\lambda$ minus that corner, so
\begin{equation}\label{eq:frecur}
f^\lambda=\sum_{c\text{ corner of }\lambda}f^{\lambda-c},\qquad f^{(1)}=1.
\end{equation}

\begin{theorem}\label{thm:standardbasis}
The polytabloids $e_T$, for $T$ standard of shape $\lambda$, are linearly independent. In particular $\dim S^\lambda\ge f^\lambda$.
\end{theorem}

\begin{proof}
Order the tabloids of shape $\lambda$ as follows. For a tabloid $t$, let $r_t(a)$ be the row containing $a$. For $t\ne t'$, let $m$ be the largest entry with $r_t(m)\ne r_{t'}(m)$, and say $t\prec t'$ if $r_t(m)<r_{t'}(m)$: at the largest
entry where they disagree, $t$ has it higher. This is a total order (it is the lexicographic order on the sequences $(r_t(n),r_t(n-1),\dots,r_t(1))$).

\emph{Claim: if $T$ is standard, then $\{T\}$ is the largest tabloid occurring in $e_T$.} The tabloids in $e_T=\sum_{\pi\in C_T}\operatorname{sgn}(\pi)\{\pi T\}$ are distinct for distinct $\pi$ (as in the proof of Theorem~\ref{thm:specht}), and
$\{T\}$ occurs with coefficient $1$. Let $\pi\in C_T$, $\pi\ne1$, and let $m$ be the largest entry whose row differs in $\{\pi T\}$ and $\{T\}$. In $\pi T$ every entry stays in its column, and an entry's position in a column is determined by its row;
so the entries larger than $m$ in the column $C$ of $m$ occupy in $\pi T$ the same cells as in $T$. As $T$ is standard, these cells are the ones below $m$. So in $\pi T$, $m$ is in one of the remaining cells of $C$, at or above its
cell in $T$, and not in the same row: strictly above. Hence $\{\pi T\}\prec\{T\}$.

Now suppose $\sum_Tc_Te_T=0$, a sum over standard $T$, with some $c_T\ne0$. A standard tableau is determined by its tabloid (sort each row), so among the $T$ with $c_T\ne0$ there is one, $T_0$, with $\{T_0\}$ largest. By the claim, $\{T_0\}$
occurs in $e_{T_0}$ with coefficient $1$, and in no other $e_T$ with $c_T\ne0$, since the tabloids of $e_T$ are $\preceq\{T\}\prec\{T_0\}$. So the coefficient of $\{T_0\}$ in the sum is $c_{T_0}\ne0$, a contradiction.
\end{proof}

The standard polytabloids also span $S^\lambda$. The classical proof straightens any $e_T$ into standard ones with the \emph{Garnir relations}; we will instead count, in the next section, that $\sum_\lambda(f^\lambda)^2=n!$. Since
$\sum_\lambda(\dim S^\lambda)^2=n!$ as well (Volume~3, Corollary~11.9), and $\dim S^\lambda\ge f^\lambda$ term by term, equality must hold everywhere.

\section{Counting standard tableaux}

\subsection*{The Robinson--Schensted correspondence}

Gilbert de Beauregard Robinson described in 1938 a way of matching permutations with pairs of tableaux, in a paper on Young's work; Craige Schensted, unaware of it, found the same algorithm in 1961 while studying increasing subsequences,
and Donald Knuth extended it in 1970 to matrices of integers. The key step is \emph{row insertion}.

Let $P$ be a \emph{partial} standard tableau: a diagram filled with distinct integers, increasing along rows and down columns. To insert $x\notin P$: if $x$ is larger than every entry of row $1$, place it at the end of row $1$ and stop. Otherwise,
let $y$ be the smallest entry of row $1$ larger than $x$; put $x$ in place of $y$ (\emph{$x$ bumps $y$}), and insert $y$ into row $2$ by the same rule, and so on, a new row being created at the bottom if needed. Write $P\leftarrow x$ for the result.

\begin{lemma}\label{lem:insertion}
$P\leftarrow x$ is a partial standard tableau, whose shape is that of $P$ plus one corner cell. Along the bumping path, the column of the bumped entry weakly decreases from each row to the next. Given $P\leftarrow x$ and the new cell,
both $P$ and $x$ can be recovered.
\end{lemma}

\begin{proof}
Suppose $x'$ enters row $i$ and bumps $y$ from column $j$; by choice of $y$, the entries of row $i$ to the left of column $j$ are smaller than $x'$, and those to the right are larger than $y$. The entry of row $i+1$ in column $j$, if any, is larger than $y$,
so when $y$ enters row $i+1$ it bumps an entry in a column $\le j$, or is appended in a column $\le j$; this proves the claim on columns. The rows stay increasing by the choice of $y$. Columns: the cell $(i,j)$ now holds $x'<y$, which
is still smaller than the entry below; above it, if $i>1$, $x'$ itself was bumped from row $i-1$ at a column $k\ge j$, and the entry of $(i-1,j)$ is now either the one that bumped $x'$ (if $k=j$) or an entry to the left of $x'$ in row $i-1$
(if $k>j$), smaller than $x'$ in both cases. The bumped $y$ lands in row $i+1$ at a column $j'\le j$; above it is $x'$ (if $j'=j$) or an entry left of column $j$ in row $i$, smaller than $x'<y$; below it is an entry larger than the
entry it replaced, which was larger than $y$. Finally the process stops by appending at the end of some row $r$, in column $\lambda_r+1$, and since columns weakly decrease along the path, $\lambda_r+1\le\lambda_{r-1}$ if the path
came from row $r-1$ (it passed there at a column $\le\lambda_{r-1}$): the new shape is a partition, and the new cell is a corner.

To recover: let $z$ be the entry of the new cell, in row $r$. If $r=1$, remove it: $x=z$. Otherwise remove it, and in row $r-1$ replace the largest entry smaller than $z$ by $z$, ejecting that entry into row $r-2$, and so on;
the entry ejected from row $1$ is $x$. This undoes insertion step by step: when $x'$ bumped $y$ in row $i$, $x'$ is the largest entry of the new row $i$ smaller than $y$, since no entry of the old row lay strictly between $x'$ and $y$.
\end{proof}

\begin{theorem}[Robinson--Schensted correspondence]\label{thm:rs}
For $\sigma\in\mathfrak S_n$, let $P(\sigma)=(\cdots((\emptyset\leftarrow\sigma(1))\leftarrow\sigma(2))\cdots)\leftarrow\sigma(n)$, and let $Q(\sigma)$ be the tableau of the same shape whose cell added at step $k$ holds $k$. Then
$\sigma\mapsto(P(\sigma),Q(\sigma))$ is a bijection from $\mathfrak S_n$ to the set of pairs of standard tableaux of the same shape. Consequently
\[
\sum_{\lambda\vdash n}(f^\lambda)^2=n!,
\]
and $\dim S^\lambda=f^\lambda$: the standard polytabloids form a basis of $S^\lambda$.
\end{theorem}

\begin{proof}
$P(\sigma)$ is standard by Lemma~\ref{lem:insertion}, with entries $1,\dots,n$. The cells of $Q(\sigma)$ are added one at a time, each a corner of the shape at that moment, so the labels increase along rows and down columns: $Q(\sigma)$
is standard. Conversely, given $(P,Q)$ of the same shape, the cell holding $n$ in $Q$ is the cell added last; Lemma~\ref{lem:insertion} recovers $\sigma(n)$ and the previous tableau, and the cell holding $n-1$ in $Q$ continues the process.
This reconstruction is defined for every pair (the reverse bumping only needs the cell to be a corner, and keeps the tableau standard by the mirror argument), and it inverts the forward map, so the map is bijective. Counting
pairs gives $\sum(f^\lambda)^2=n!$, and the dimension count after Theorem~\ref{thm:standardbasis} finishes the proof.
\end{proof}

For example, $\sigma=31452$ (in one-line notation) gives, step by step, $P$ with rows $3$; then $1/3$ (the $1$ bumps the $3$); then $14/3$; then $145/3$; then $125/34$ (the $2$ bumps the $4$, which goes to the end of row $2$). The cells were
added in the order $(1,1),(2,1),(1,2),(1,3),(2,2)$, so $Q$ has rows $134/25$.

\begin{theorem}[Schensted, 1961]\label{thm:schensted}
The length of the first row of $P(\sigma)$ is the length of the longest increasing subsequence of $\sigma(1),\sigma(2),\dots,\sigma(n)$.
\end{theorem}

\begin{proof}
We show by induction on $k$ that after inserting $\sigma(1),\dots,\sigma(k)$, the first row $r_1<r_2<\dots<r_\ell$ satisfies: increasing subsequences of $\sigma(1),\dots,\sigma(k)$ have length at most $\ell$, and for each $j\le\ell$, $r_j$ is the smallest
possible last term of an increasing subsequence of length $j$. This is clear for $k=1$. Insert $x=\sigma(k+1)$, and let $L=1+\#\{j:r_j<x\}$. An increasing subsequence ending at $x$ of length $m$ consists of one of length $m-1$
ending below $x$, and then its last term is $\ge r_{m-1}$, so $r_{m-1}<x$ and $m\le L$; conversely, a subsequence of length $L-1$ ending at $r_{L-1}<x$ extends by $x$. So the longest increasing subsequence ending at $x$ has length exactly $L$.
Row insertion puts $x$ in column $L$ of row $1$, replacing $r_L>x$ or appending if $L=\ell+1$. The minimal last term for length $L$ becomes $x$; for $m<L$, $x>r_{L-1}\ge r_m$ does not improve $r_m$; and for $m>L$ no subsequence of
length $m$ ends at $x$. The invariant holds after step $k+1$.
\end{proof}

For $31452$, the first row of $P$ is $125$ and the longest increasing subsequences, $345$ and $145$, have length $3$. Note that the first row of $P$ is generally \emph{not} an increasing subsequence of $\sigma$: $1,2,5$ do not appear in that
order. Schensted also proved that the number of rows is the length of the longest \emph{decreasing} subsequence; the Erdős--Szekeres theorem follows, and Exercise~\ref{ex:patience} gives a direct route.

\subsection*{The hook length formula}

\begin{definition}\label{def:hook}
The \emph{hook} of a cell $(i,j)$ of $\lambda$ is the cell itself, the cells to its right in row $i$ (its \emph{arm}) and the cells below it in column $j$ (its \emph{leg}). Its \emph{hook length} is
$h_{ij}=(\lambda_i-j)+(\lambda'_j-i)+1$.
\end{definition}

\begin{figure}[ht]
\centering
\fig{ch09-hooks}
\caption{Left: the diagram of $\lambda=(5,4,2,1)$ with the hook length in each cell; the hook of the cell $(1,2)$ is shaded, of length $6$. By the hook length formula $f^\lambda=12!/(8\cdot6\cdot4\cdot3\cdot1\cdot6\cdot4\cdot2\cdot1\cdot3\cdot1\cdot1)=5775$.
Right: one run of the hook walk of Theorem~\ref{thm:hook}, from $(1,1)$ to $(3,1)$ (down its leg), then to the corner $(3,2)$ (along its arm). The four corners are marked with dots.}
\label{fig:hooks}
\end{figure}

\begin{theorem}[Hook length formula; Frame, Robinson and Thrall, 1954]\label{thm:hook}
For $\lambda\vdash n$,
\[
f^\lambda=\frac{n!}{\prod_{(i,j)\in\lambda}h_{ij}}.
\]
\end{theorem}

The proof below is due to Curtis Greene, Albert Nijenhuis and Herbert Wilf (1979), and it is probabilistic. Write $F(\lambda)$ for the right-hand side. By \eqref{eq:frecur} and induction on $n$ (with $F((1))=1$), it suffices to show that
$F$ satisfies the same recursion, that is,
\begin{equation}\label{eq:hooksum}
\sum_{c\text{ corner}}\frac{F(\lambda-c)}{F(\lambda)}=1.
\end{equation}
The idea is to find a random process that ends at a corner $c$ with probability exactly $F(\lambda-c)/F(\lambda)$; the probabilities then add up to $1$ for free.

\emph{The hook walk.} Choose a cell of $\lambda$ uniformly at random, among $n$. From the current cell, if it is not a corner, move to a cell chosen uniformly among the $h-1$ other cells of its hook. Each move goes strictly right or strictly
down, so the walk stops, at a corner (the corners are exactly the cells with $h=1$).

\begin{proof}[Proof of Theorem~\ref{thm:hook}]
\emph{The target.} Let $c=(\alpha,\beta)$ be a corner. Removing $c$ changes the hook lengths only of the cells $(i,\beta)$, $i<\alpha$, and $(\alpha,j)$, $j<\beta$, each by $-1$, and removes the factor $h_{\alpha\beta}=1$. So
\[
\frac{F(\lambda-c)}{F(\lambda)}=\frac1n\prod_{i<\alpha}\frac{h_{i\beta}}{h_{i\beta}-1}\prod_{j<\beta}\frac{h_{\alpha j}}{h_{\alpha j}-1}
=\frac1n\prod_{i<\alpha}\Bigl(1+\frac1{h_{i\beta}-1}\Bigr)\prod_{j<\beta}\Bigl(1+\frac1{h_{\alpha j}-1}\Bigr).
\]
Expanding the products,
\begin{equation}\label{eq:hooktarget}
\frac{F(\lambda-c)}{F(\lambda)}=\frac1n\sum_{A,B}\ \prod_{i\in A}\frac1{h_{i\beta}-1}\prod_{j\in B}\frac1{h_{\alpha j}-1},
\end{equation}
over all subsets $A\subseteq\{1,\dots,\alpha-1\}$ and $B\subseteq\{1,\dots,\beta-1\}$.

\emph{The walk.} For a run of the walk ending at $c$, let $\bar A$ be the set of rows it visits and $\bar B$ the set of columns; $\alpha\in\bar A$, $\beta\in\bar B$, and the start is $(\min\bar A,\min\bar B)$. We claim that for given $\bar A\ni\alpha$,
$\bar B\ni\beta$, the probability that the walk started at $(a,b)=(\min\bar A,\min\bar B)$ ends at $c$ with these sets of rows and columns is
\[
p(\bar A,\bar B)=\prod_{i\in\bar A\setminus\{\alpha\}}\frac1{h_{i\beta}-1}\prod_{j\in\bar B\setminus\{\beta\}}\frac1{h_{\alpha j}-1}.
\]
By induction on $|\bar A|+|\bar B|$. If both are singletons, the walk starts at the corner and stops: probability $1$. Otherwise the first move, of probability $1/(h_{ab}-1)$, goes down to the next row $a_2$ of $\bar A$ in column $b$,
or right to the next column $b_2$ of $\bar B$ in row $a$ (only one of these is possible if $a=\alpha$ or $b=\beta$). By induction,
\[
\frac{p(\bar A\setminus\{a\},\bar B)+p(\bar A,\bar B\setminus\{b\})}{h_{ab}-1}
=p(\bar A,\bar B)\,\frac{(h_{a\beta}-1)+(h_{\alpha b}-1)}{h_{ab}-1},
\]
where a missing term corresponds to $a=\alpha$ or $b=\beta$, and then $h_{\alpha\beta}-1=0$ accounts for it. Since $c$ is a corner, $\lambda_\alpha=\beta$ and $\lambda'_\beta=\alpha$, and
\[
(h_{a\beta}-1)+(h_{\alpha b}-1)=(\lambda_a-\beta+\alpha-a)+(\beta-b+\lambda'_b-\alpha)=\lambda_a-b+\lambda'_b-a=h_{ab}-1.
\]
So the fraction is $1$, which proves the claim.

\emph{Conclusion.} The start is uniform, so the probability of ending at $c$ is $\frac1n\sum p(\bar A,\bar B)$ over all $\bar A\ni\alpha$ with $\max\bar A=\alpha$ and $\bar B\ni\beta$ with $\max\bar B=\beta$. Writing $A=\bar A\setminus\{\alpha\}$ and
$B=\bar B\setminus\{\beta\}$, this is exactly \eqref{eq:hooktarget}. The walk ends at some corner, so \eqref{eq:hooksum} holds, and $f^\lambda=F(\lambda)$ by induction.
\end{proof}

For $\lambda=(2,2)$ the hooks are $3,2,2,1$ and $f=24/12=2$; for the staircase $(3,2,1)$ they are $5,3,1,3,1,1$ and $f=720/45=16$. Exercise~\ref{ex:frobeniusdim} derives from the theorem the formula with which Frobenius computed dimensions in 1900.

\section{The character of a transposition}

Let $\chi^\lambda$ be the character of $S^\lambda$. Computing the whole character table of $\mathfrak S_n$ is a story of its own (the Murnaghan--Nakayama rule); for shuffling cards we need one column, the value on a transposition. There is a beautiful formula.
The \emph{content} of a cell $(i,j)$ is $j-i$: $0$ on the diagonal, positive above, negative below.

\begin{theorem}[Frobenius, 1900]\label{thm:transposition}
For a transposition $\tau$ and $\lambda\vdash n$,
\[
\frac{\chi^\lambda(\tau)}{f^\lambda}=\frac1{\binom n2}\Bigl(\sum_i\binom{\lambda_i}2-\sum_j\binom{\lambda'_j}2\Bigr)=\frac1{\binom n2}\sum_{(i,j)\in\lambda}(j-i).
\]
\end{theorem}

\begin{proof}
Let $X\in\C[\mathfrak S_n]$ be the sum of all $\binom n2$ transpositions. It is a class sum, hence central, and by Schur's lemma it acts on the irreducible module $S^\lambda$ by a scalar $s$. Taking traces, $sf^\lambda=\binom n2\chi^\lambda(\tau)$,
since all transpositions are conjugate. So it suffices to show that $s$ is the bracket.

Fix a tableau $T$ of shape $\lambda$, and compare the coefficients of the tabloid $\{T\}$ on both sides of $Xe_T=s\,e_T$. On the right it is $s$, as $\{T\}$ has coefficient $1$ in $e_T$. On the left,
\[
Xe_T=\sum_t\sum_{\pi\in C_T}\operatorname{sgn}(\pi)\{t\pi T\},
\]
and the pair $(t,\pi)$ contributes to $\{T\}$ if and only if $t\pi\in R_T$. We sort the transpositions $t=(a\,b)$ into three kinds.
\begin{itemize}
\item $a,b$ in the same row of $T$: then $t\pi\in R_T$ means $\pi\in tR_T=R_T$, so $\pi\in C_T\cap R_T=\{1\}$, contributing $+1$.
\item $a,b$ in the same column: then $t\in C_T$, and $t\pi\in R_T\cap C_T=\{1\}$ forces $\pi=t$, contributing $\operatorname{sgn}(t)=-1$.
\item $a,b$ in different rows and different columns: suppose $t\pi\in R_T$, that is, $\{t\pi T\}=\{T\}$. For an entry $e\notin\{a,b\}$, $t$ does not move it, so $e$ has the same cell in $t\pi T$ as in $\pi T$: in its own column, since $\pi\in C_T$,
and in its own row, since $\{t\pi T\}=\{T\}$; that is, in its own cell, so $\pi$ fixes $e$. Hence $\pi\in\{1,t\}\cap C_T=\{1\}$, and then $t\in R_T$, which is false. No contribution.
\end{itemize}
So $s$ is the number of transpositions inside rows minus the number inside columns, $\sum_i\binom{\lambda_i}2-\sum_j\binom{\lambda'_j}2$. Finally, write the content as $j-i=(j-1)-(i-1)$ and sum over all cells: row by row, $\sum_{j\le\lambda_i}(j-1)=\binom{\lambda_i}2$, and column by column, $\sum_{i\le\lambda'_j}(i-1)=\binom{\lambda'_j}2$.
\end{proof}

Checks: for $\lambda=(n)$ the ratio is $1$ (trivial representation), for $(1^n)$ it is $-1$ (sign). For $(n-1,1)$ the contents are $0,1,\dots,n-2$ and $-1$, the ratio is $\bigl(\binom{n-1}2-1\bigr)/\binom n2=\frac{n-3}{n-1}$, and indeed the standard
representation has $\chi(\tau)=(\text{fixed points})-1=n-3$ and dimension $n-1$.

\begin{corollary}\label{cor:conjugate}
$\chi^{\lambda'}(\tau)/f^{\lambda'}=-\chi^\lambda(\tau)/f^\lambda$, and $f^{\lambda'}=f^\lambda$.
\end{corollary}

\begin{proof}
Transposing the diagram swaps rows and columns, so it negates contents; and it maps standard tableaux of shape $\lambda$ bijectively to those of shape $\lambda'$.
\end{proof}

In fact $S^{\lambda'}\cong S^\lambda\otimes\operatorname{sgn}$, but we will not need it. What Volume~11 uses is the eigenvalue picture. Random transpositions act on $\C[\mathfrak S_n]$ by multiplication by $\frac1{\binom n2}X$ (or a lazy version of it),
and by Wedderburn (Corollary~\ref{cor:wedderburnCG}), $\C[\mathfrak S_n]$ splits into $f^\lambda$ copies of each $S^\lambda$, on which $X/\binom n2$ is the scalar $\chi^\lambda(\tau)/f^\lambda$. So the transition matrix of the shuffle is diagonalised by
representation theory, with eigenvalue $\chi^\lambda(\tau)/f^\lambda$ of multiplicity $(f^\lambda)^2$, and everything is now explicit.

\begin{beauty}
In 1961 Stanislaw Ulam asked how long the longest increasing subsequence $L_n$ of a random permutation of $\{1,\dots,n\}$ is, and simulations suggested about $2\sqrt n$. By Theorem~\ref{thm:rs} and Theorem~\ref{thm:schensted}, $L_n$
is the first row of a random Young diagram drawn with probability $(f^\lambda)^2/n!$, the \emph{Plancherel measure}. In 1977 Benjamin Logan and Larry Shepp, and independently Anatoly Vershik and Sergei Kerov, proved that, scaled by
$\sqrt n$, these diagrams converge to a fixed curve (rotated by $45^\circ$, the graph of $\frac2\pi\bigl(u\arcsin\frac u2+\sqrt{4-u^2}\bigr)$ for $|u|\le2$), and that $L_n/\sqrt n\to2$. The cover of this book shows a random Plancherel diagram. Then in
1999 Jinho Baik, Percy Deift and Kurt Johansson proved that $(L_n-2\sqrt n)/n^{1/6}$ converges in law to the Tracy--Widom distribution, which describes the largest eigenvalue of a large random Hermitian matrix. Sorting a deck of
cards and the spectrum of random matrices obey the same law; random unitary matrices return in Chapter~11.
\end{beauty}

\section{Exercises}

\begin{exercise}[$\star$]\label{ex:catalan}
Show with the hook length formula that $f^{(n,n)}=\frac1{n+1}\binom{2n}n$, the Catalan number. Find a direct bijection between standard tableaux of shape $(n,n)$ and sequences of $n$ opening and $n$ closing brackets that are correctly matched.
\end{exercise}

\begin{exercise}[$\star$]\label{ex:smalltables}
List the partitions of $4$ and compute $f^\lambda$ and $\chi^\lambda(\tau)$ for each with the hook length formula and Theorem~\ref{thm:transposition}. Match them with the rows of the character table of $\mathfrak S_4$ (Volume~3, Example~11.13). Then compute
the dimensions of the irreducible representations of $\mathfrak S_5$ and check that their squares add up to $120$.
\end{exercise}

\begin{exercise}[$\star$]\label{ex:standardrep}
Show that $M^{(n-1,1)}\cong S^{(n)}\oplus S^{(n-1,1)}$, using only the remark after Theorem~\ref{thm:specht} and dimensions. Then show with the hook length formula that $f^{(n-j,j)}=\binom nj-\binom n{j-1}$ for $0\le j\le n/2$.
\end{exercise}

\begin{exercise}[$\star\star$]\label{ex:youngmult}
Let $\mu\vdash n$. Show that every module map $\theta:S^\mu\to M^\mu$ is a multiple of the inclusion, and deduce that $S^\mu$ occurs exactly once in $M^\mu$ and that every other irreducible constituent $S^\lambda$ of $M^\mu$ has $\lambda\triangleright\mu$.
\end{exercise}

\begin{exercise}[$\star\star$]\label{ex:patience}
\emph{Patience sorting.} Deal the values $\sigma(1),\dots,\sigma(n)$ one at a time onto piles: each card goes on the leftmost pile whose top card is larger, or starts a new pile on the right. Show that each pile read from bottom to top is a
decreasing subsequence, that the number of piles is the length of the longest increasing subsequence, and that the tops of the piles are always the first row of the Robinson--Schensted tableau $P$. Deduce the Erdős--Szekeres theorem (1935):
every sequence of $rs+1$ distinct numbers has an increasing subsequence of length $r+1$ or a decreasing one of length $s+1$.
\end{exercise}

\begin{exercise}[$\star\star\star$]\label{ex:frobeniusdim}
Let $\lambda=(\lambda_1,\dots,\lambda_k)$ and $l_i=\lambda_i+k-i$, so $l_1>l_2>\dots>l_k\ge0$. Show that the hook lengths in row $i$ are the numbers $1,2,\dots,l_i$ with the $k-i$ numbers $l_i-l_j$, $j>i$, removed. Deduce Frobenius' formula
\[
f^\lambda=\frac{n!\prod_{i<j}(l_i-l_j)}{l_1!\,l_2!\cdots l_k!}.
\]
\end{exercise}

\section{Hints}

\noindent\textbf{\ref{ex:catalan}.} The hooks in row $1$ are $n+1,\dots,2$, in row $2$ they are $n,\dots,1$. Read the tableau's entries in increasing order.

\noindent\textbf{\ref{ex:smalltables}.} The contents of $(3,1)$ are $0,1,2,-1$. For $\mathfrak S_5$, there are $7$ partitions.

\noindent\textbf{\ref{ex:standardrep}.} Which partitions dominate $(n-1,1)$? For the hooks, treat the first $j$ columns separately.

\noindent\textbf{\ref{ex:youngmult}.} By Lemma~\ref{lem:kappa}(b), $\theta(e_T)=\kappa_T\Theta(\{T\})$ is a multiple of $e_T$.

\noindent\textbf{\ref{ex:patience}.} Compare the pile tops with the invariant in the proof of Theorem~\ref{thm:schensted}.

\noindent\textbf{\ref{ex:frobeniusdim}.} The hook lengths of a row decrease strictly; show that $h_{ic}\ne l_i-l_j$ by comparing $\lambda_j$ with $c$, then count.

\section{Solutions}

\begin{solution}{ex:catalan}
\idea Two rows of hooks, then a bracket for each entry according to its row.
\writeup The cell $(1,j)$ has arm $n-j$ and leg $1$, so hook $n-j+2$; the cell $(2,j)$ has hook $n-j+1$. The product is $(n+1)!\,n!$, and $f^{(n,n)}=\frac{(2n)!}{(n+1)!\,n!}=\frac1{n+1}\binom{2n}n$. For the bijection, read $1,2,\dots,2n$ and write an opening bracket
for each entry of row $1$, a closing one for each entry of row $2$. The tableau is standard if and only if the rows increase (automatic once we choose which numbers go in which row and sort them) and each column increases, that is, the $j$-th
entry of row $2$ is larger than the $j$-th entry of row $1$: at every moment, at least as many entries of row $1$ have been read as of row $2$. This is exactly the condition for brackets to match, and the inverse map puts the positions of
the opening brackets in row $1$.
\end{solution}

\begin{solution}{ex:smalltables}
\idea Hooks give dimensions, contents give the transposition column; Corollary~\ref{cor:conjugate} halves the work.
\writeup For $n=4$, $\binom42=6$:
\begin{center}
\begin{tabular}{@{}l@{\enspace}c@{\enspace}c@{\enspace}c@{\enspace}c@{\enspace}c@{}}
$\lambda$ & $(4)$ & $(3,1)$ & $(2,2)$ & $(2,1,1)$ & $(1^4)$\\\hline
$f^\lambda$ & $1$ & $3$ & $2$ & $3$ & $1$\\
contents & $6$ & $2$ & $0$ & $-2$ & $-6$\\
$\chi^\lambda(\tau)$ & $1$ & $1$ & $0$ & $-1$ & $-1$
\end{tabular}
\end{center}
(for $(3,1)$ the hooks are $4,2,1,1$, so $f=24/8=3$, and $\chi(\tau)=3\cdot2/6=1$). These match the trivial character, $\chi_3$ (the tetrahedron), the character $(2,0,2,-1,0)$ coming from $\mathfrak S_3$, $\chi_3'=\varepsilon\chi_3$ and $\varepsilon$:
the values at $(12)$ in Volume~3 are $1,1,0,-1,-1$, and the dimensions separate the rest. For $n=5$ the partitions $(5),(4,1),(3,2),(3,1,1),(2,2,1),(2,1,1,1),(1^5)$ give $f=1,4,5,6,5,4,1$ (for instance $(3,2)$ has hooks $4,3,1,2,1$ and $f=120/24=5$; $(3,1,1)$ has hooks $5,2,1,2,1$ and
$f=120/20=6$), and $1+16+25+36+25+16+1=120$.
\end{solution}

\begin{solution}{ex:standardrep}
\idea Dominance leaves only two candidates, and the dimension decides the multiplicities.
\writeup The only partitions dominating $(n-1,1)$ are $(n)$ and $(n-1,1)$ itself. So $M^{(n-1,1)}\cong m\,S^{(n)}\oplus S^{(n-1,1)}$ for some $m\ge0$, with $S^{(n-1,1)}$ once (the remark after Theorem~\ref{thm:specht}). Dimensions give $n=m+(n-1)$, so $m=1$;
concretely, $S^{(n)}$ is the line spanned by $\sum_a\mathbf e_a$ and $S^{(n-1,1)}$ the hyperplane $\sum x_a=0$. For the second part, in the diagram of $(n-j,j)$ the cell $(1,c)$ has hook $n-j-c+2$ for $c\le j$ and $n-j-c+1$ for $c>j$, and the cell $(2,c)$ has hook $j-c+1$.
So the product of the hooks is
\[
\frac{(n-j+1)!}{n-2j+1}\cdot j!,\qquad\text{and}\qquad f^{(n-j,j)}=\frac{n!\,(n-2j+1)}{j!\,(n-j+1)!}=\binom nj-\binom n{j-1},
\]
the last equality because $\binom nj-\binom n{j-1}=\frac{n!}{j!\,(n-j+1)!}\bigl((n-j+1)-j\bigr)$. (In row $1$, the hooks are $n-j+1,\dots,n-2j+2$ for the first $j$ cells and $n-2j,\dots,1$ for the others: all of $1,\dots,n-j+1$ except $n-2j+1$.)
\end{solution}

\begin{solution}{ex:youngmult}
\idea A map out of $S^\mu$ is determined by where $e_T$ goes, and $\kappa_T$ squeezes the image onto the line through $e_T$.
\writeup Extend $\theta$ by zero on $(S^\mu)^\perp$ to $\Theta:M^\mu\to M^\mu$. Then $\theta(e_T)=\kappa_T\Theta(\{T\})=c\,e_T$ by Lemma~\ref{lem:kappa}(b). For any $\sigma$, $\theta(e_{\sigma T})=\sigma\theta(e_T)=c\,e_{\sigma T}$, so the same $c$ works for all
polytabloids, and $\theta=c\cdot\mathrm{incl}$. Hence $\dim\operatorname{Hom}(S^\mu,M^\mu)=1$, which is the multiplicity of $S^\mu$ in $M^\mu$ (Schur's lemma, decomposing $M^\mu$ into irreducibles). If $S^\lambda$ is a constituent, there is a non-zero
map $S^\lambda\to M^\mu$, and the proof of Theorem~\ref{thm:specht} gives $\lambda\trianglerighteq\mu$; for $\lambda\ne\mu$, $\lambda\triangleright\mu$.
\end{solution}

\begin{solution}{ex:patience}
\idea The pile tops are the minimal last terms of the proof of Schensted's theorem.
\writeup A card goes on a pile only if the top is larger, so piles decrease from bottom to top; the cards of a pile were dealt in that order, so a pile is a decreasing subsequence. The tops increase from left to right: a new pile is started only
when the card is larger than all tops, and when $x$ lands on the leftmost pile with top $>x$, the pile to its left has top $<x$ and the one to its right had top $>$ the old top $>x$. So the tops form an increasing row, and placing $x$ on the leftmost
top larger than $x$ is exactly the first-row step of row insertion: the tops are the first row of $P$, by induction. By Theorem~\ref{thm:schensted}, the number of piles is the longest increasing subsequence. Finally, if a sequence of $rs+1$ distinct numbers
has no increasing subsequence of length $r+1$, there are at most $r$ piles, so some pile has at least $s+1$ cards: a decreasing subsequence of length $s+1$.
\end{solution}

\begin{solution}{ex:frobeniusdim}
\idea There are exactly $l_i$ numbers in play; show that the two lists do not overlap.
\writeup The hook lengths in row $i$ strictly decrease from left to right, from $h_{i1}=\lambda_i-1+\lambda'_1-i+1=\lambda_i+k-i=l_i$ (as $\lambda'_1=k$) down to $h_{i\lambda_i}\ge1$: $\lambda_i$ distinct numbers in $\{1,\dots,l_i\}$. The $k-i$ numbers $l_i-l_j$, $j>i$, are distinct
and also in $\{1,\dots,l_i\}$. As $\lambda_i+(k-i)=l_i$, it suffices to show that $h_{ic}\ne l_i-l_j$ for every $c\le\lambda_i$ and $j>i$. Now
\[
h_{ic}-(l_i-l_j)=(\lambda_i-c+\lambda'_c-i+1)-(\lambda_i-\lambda_j+j-i)=\lambda_j+\lambda'_c+1-c-j.
\]
If $\lambda_j\ge c$, the cell $(j,c)$ exists, so $\lambda'_c\ge j$, and the difference is $\ge1$. If $\lambda_j\le c-1$, the cell $(j,c)$ does not exist, so $\lambda'_c\le j-1$, and the difference is $\le-1$. So the hook lengths in row $i$ multiply to
$l_i!/\prod_{j>i}(l_i-l_j)$, and over all rows $\prod h=\prod_il_i!/\prod_{i<j}(l_i-l_j)$. The hook length formula gives Frobenius' formula.
\end{solution}

\section*{Sources and further reading}
\addcontentsline{toc}{section}{Sources and further reading}

A.~Young, ``On quantitative substitutional analysis'', \emph{Proc.\ London Math.\ Soc.} 33 (1901), 97--146, and 34 (1902), 361--397, the first two papers of nine; the biographical facts and the quotation of Turnbull are from the MacTutor
biography of Alfred Young (J.~J.~O'Connor and E.~F.~Robertson). G.~Frobenius, ``Über die Charaktere der symmetrischen Gruppe'', \emph{Sitzungsber.\ Preuss.\ Akad.\ Wiss.} (1900), 516--534; W.~Specht, ``Die irreduziblen Darstellungen der
symmetrischen Gruppe'', \emph{Math.\ Z.} 39 (1935), 696--711; G.~D.~James, \emph{The Representation Theory of the Symmetric Groups}, Lecture Notes in Mathematics 682 (Springer, 1978). G.~de~B.~Robinson, ``On the representations of the symmetric
group'', \emph{Amer.\ J.\ Math.} 60 (1938), 745--760; C.~Schensted, ``Longest increasing and decreasing subsequences'', \emph{Canad.\ J.\ Math.} 13 (1961), 179--191; D.~E.~Knuth, ``Permutations, matrices, and generalized Young tableaux'',
\emph{Pacific J.\ Math.} 34 (1970), 709--727. J.~S.~Frame, G.~de~B.~Robinson and R.~M.~Thrall, ``The hook graphs of the symmetric group'', \emph{Canad.\ J.\ Math.} 6 (1954), 316--324; C.~Greene, A.~Nijenhuis and H.~S.~Wilf, ``A probabilistic proof of a
formula for the number of Young tableaux of a given shape'', \emph{Adv.\ Math.} 31 (1979), 104--109. B.~F.~Logan and L.~A.~Shepp, ``A variational problem for random Young tableaux'', \emph{Adv.\ Math.} 26 (1977), 206--222; A.~M.~Vershik and S.~V.~Kerov,
``Asymptotics of the Plancherel measure of the symmetric group and the limiting form of Young tableaux'', \emph{Soviet Math.\ Dokl.} 18 (1977), 527--531; J.~Baik, P.~Deift and K.~Johansson, ``On the distribution of the length of the longest increasing
subsequence of random permutations'', \emph{J.\ Amer.\ Math.\ Soc.} 12 (1999), 1119--1178.

For more: B.~E.~Sagan, \emph{The Symmetric Group} (Springer, 2nd edition 2001), Chapters~2 and~3, whose route we followed; W.~Fulton, \emph{Young Tableaux} (Cambridge University Press, 1997); D.~Romik, \emph{The Surprising Mathematics of Longest
Increasing Subsequences} (Cambridge University Press, 2015), for the beauty box.
